% !PRINTED 614
\weberpage{614}
%<<CONT
Es sei $\mathfrak{m}$ ein ganzes Ideal in $\Omega$, und $A$ bestehe aus allen ganzen und gebrochenen Zahlen $\alpha$ in $\Omega$, die zu $\mathfrak{m}$ teilerfremd sind und der Bedingung genügen:
\weq{1}{\alpha \equiv 1 \quad (\mathrm{mod}\ \mathfrak{m}).}

Nehmen wir die Primteiler von $\mathfrak{m}$ in den Exkludenten, und bezeichnen mit $O$ die Gruppe der Zahlen in $\Omega$, so zerfällt $O$ nach dem Modul $\mathfrak{m}$ in $\psi(\mathfrak{m})$ Zahlklassen, worin $\psi(\mathfrak{m})$ die Bedeutung wie in Bd.\ II, \S\,168 hat, nämlich, wenn $\mathfrak{p}$ die Primteiler von $\mathfrak{m}$ durchläuft:
\weq{2}{\psi(\mathfrak{m}) = N(\mathfrak{m})\ \Pi\Bigl(1 - \frac{1}{N(\mathfrak{p})}\Bigr).}

In die Gruppe $A$ gehören nicht bloß einzelne Zahlen, sondern Zahlklassen nach dem Modul $\mathfrak{m}$. Hiernach ist:
\weq{3}{(O, A) = \psi(\mathfrak{m}).}

Um die Klassenzahl für unsere Gruppe zu bestimmen, wenden wir die Sätze des \S\,161 an.

Danach ist, wenn $\overline{E}$ die Gruppe der funktionalen Einheiten und $E$ die Gruppe der numerischen Einheiten bedeutet:
\[
(\overline{O}, \overline{E}A) = (\overline{O}, \overline{E}O)\,(\overline{E}O, \overline{E}A),
\]
\[
(\overline{E}O, \overline{E}A) = (\overline{E}O, \overline{E}EA) = (O, EA).
\]

Weiter ist
\[
(O, EA)\,(EA, A) = (O, A) = \psi(\mathfrak{m}).
\]
$(EA, A) = e$ ist die Zahl der nach $\mathfrak{m}$ inkongruenten Einheiten, und
\[
(\overline{O}, \overline{E}O) = H
\]
die Klassenzahl des Körpers $\Omega$,
\[
(\overline{O}, \overline{E}A) = h
\]
die Klassenzahl nach $A$. Daraus ergibt sich:
\weq{4}{h = \frac{H\,\psi(\mathfrak{m})}{e}.}

Die Zahl $e$ ist im allgemeinen gleich der Anzahl der Einheiten in $\Omega$, also $= 2$, und in den beiden Ausnahmefällen $\Delta = -4$, $\Delta = -3$ ist $e = 4$ und $e = 6$. Die Zahl $e$ ist kleiner, wenn unter den Einheiten Kongruente (modulo $\mathfrak{m}$) vorkommen, was nur möglich ist, wenn $\mathfrak{m}$ ein Teiler von 2 oder von 3 ist.
